\documentclass{article}
\usepackage[utf8]{inputenc}
\usepackage{a4wide}

\begin{document}

\section*{เล็ก กลาง ใหญ่}

นิพพิชฌน์เปิดร้านขายลูกชิ้นหมูปิ้งเสียบไม้ โดยลูกชิ้นหมูปิ้ง 1 ไม้จะประกอบด้วยลูกชิ้นสามลูก:

ลูกใหญ่, ลูกกลาง, และลูกเล็ก ซึ่งลูกชิ้นแต่ละลูกจะต้องมีรัศมีที่ต่างกัน นิพพิชฌน์มีลูกชิ้นจำนวน n

ลูก ซึ่งแต่ละลูกมีรัศมี r1, r2, ..., rn



\ul{ตัวอย่างเช่น}หากเธอมีลูกชิ้นรัศมี 1, 2, 3 เธอจะทำลูกชิ้นหมูปิ้งได้ 1 ไม้ แต่ถ้ามีลูกชิ้นรัศมี

2, 2, 3 หรือ 2, 2, 2 จะไม่สามารถทำลูกชิ้นหมูปิ้งได้



\ul{งานของคุณ}: ช่วยนิพพิชฌน์คำนวณว่าจากลูกชิ้นที่มีอยู่สามารถทำลูกชิ้นหมูปิ้งได้มากที่สุดกี่ไม้

โดยแต่ละลูกชิ้นแต่ละไม้ต้องมีรัศมีของลูกชิ้นสามลูกที่แตกต่างกัน



\subsection*{Input}
บรรทัดแรก: จำนวนเต็ม n แสดงจำนวนลูกชิ้นที่มี โดยที่ 1 ≤ n ≤ 10\^{}5



บรรทัดที่สอง: r1 r2 \ldots{} rn เป็นจำนวนเต็มแสดงรัศมีของลูกชิ้นแต่ละลูก โดยที่

1 ≤ r{[}i{]} ≤ 10\^{}9



\subsection*{Output}
บรรทัดแรก: k แสดงจำนวนไม้ของลูกชิ้นหมูปิ้ง



บรรทัดต่อๆ ไปจำนวน k บรรทัด: รัศมีของลูกชิ้นในแต่ละไม้ เรียงจาก ใหญ่ไปเล็ก



\end{document}
